\chapter{Metodología para el diseño de COBRIA}
\label{ch:metodologia-diseno}

\section{Enfoque de ciencia del diseño}

COBRIA se estudia como un artefacto de ingeniería y no únicamente como una descripción
de prácticas observadas. La ciencia del diseño resulta apropiada cuando la investigación
busca construir y evaluar un artefacto destinado a resolver un problema relevante. Ese
artefacto puede adoptar la forma de constructos, modelos, métodos o instanciaciones
\parencite{hevner2004designscience}. COBRIA comprende las cuatro formas: propone un
vocabulario, formaliza relaciones, prescribe un método de diseño y se instancia en
sistemas ejecutables.

El proceso sigue seis actividades adaptadas de la metodología de
\textcite{peffers2007dsrm}:

\begin{enumerate}[leftmargin=*]
  \item identificar el problema de autoridad, duplicación, consulta y trazabilidad;
  \item definir objetivos verificables para una solución;
  \item diseñar y desarrollar los componentes de COBRIA;
  \item demostrar su aplicación en dos casos anonimizados;
  \item evaluar beneficios, costos y fallos mediante experimentos;
  \item comunicar resultados, límites y material de replicación.
\end{enumerate}

El desarrollo no es estrictamente lineal. Los hallazgos de evaluación pueden exigir
reformular un contrato, retirar una proyección o reducir una abstracción. Este ciclo
evita presentar como principio general una decisión que solo funciona en un caso.

\section{Proceso de abstracción desde los casos}

Los dos sistemas de estudio contienen implementaciones con diferentes grados de
uniformidad. El Sistema A combina componentes heredados y especializados; el Sistema B
utiliza metadatos y generación con mayor regularidad. COBRIA no copia literalmente
ninguna de las dos estructuras. Su derivación sigue cuatro pasos:

\begin{enumerate}[leftmargin=*]
  \item \textbf{Inventario:} identificar entidades, accesos, servicios, operaciones,
  índices, catálogos y procesos asíncronos.
  \item \textbf{Normalización conceptual:} traducir nombres internos a términos
  reconocibles, como Entidad, Repositorio, Mapeador de datos, Capa de servicios y Proyección.
  \item \textbf{Comparación:} separar mecanismos comunes de decisiones exclusivas de
  un dominio o de deuda histórica.
  \item \textbf{Generalización cautelosa:} conservar únicamente propiedades que puedan
  definirse, implementarse y posteriormente evaluarse.
\end{enumerate}

La correspondencia terminológica inicial se muestra en la
\cref{tab:correspondencia-componentes}. El término interno \emph{Dater}, útil para
describir los casos, se sustituye en la propuesta por Repositorio/Mapeador de datos, una
terminología más próxima a la literatura de arquitectura empresarial
\parencite{fowler2002poeaa}.

\begin{table}[H]
\centering
\caption{Normalización conceptual de componentes observados.}
\label{tab:correspondencia-componentes}
\begin{tabularx}{\textwidth}{p{3.2cm}p{4.2cm}X}
\toprule
Nombre observado & Término COBRIA & Responsabilidad principal \\
\midrule
Objeto & Canónico Entidad & Representar identidad, estado semántico, versión y ciclo de vida.\\
Dater & Repositorio/Mapeador de datos & Traducir representaciones y encapsular persistencia y consulta.\\
Servicio & Dominio/Aplicación Servicio & Aplicar políticas, componer datos y coordinar capacidades.\\
CasoDeUso & Caso de uso de aplicación & Ejecutar una intención autorizada y completa.\\
Servicio indexador & Gestor de proyecciones & Mantener, verificar y reconstruir proyecciones.\\
thingsData & Registro de catálogos & Administrar vocabularios finitos y datos de referencia.\\
DataConnect & Pasarela de almacenamiento & Adaptar el motor físico y sus primitivas.\\
\bottomrule
\end{tabularx}
\end{table}

\section{Criterios de diseño}

Cada componente propuesto debe satisfacer cinco criterios:

\begin{description}[leftmargin=3.4cm,style=nextline]
  \item[Comprensibilidad] Su propósito debe poder explicarse sin conocer una ruta física.
  \item[Verificabilidad] Deben existir invariantes o resultados observables.
  \item[Reemplazabilidad] El dominio no debe depender innecesariamente de un proveedor.
  \item[Proporcionalidad] La complejidad añadida debe responder a una necesidad medible.
  \item[Trazabilidad] Debe conocerse el origen y la versión de toda derivación relevante.
\end{description}

Estos criterios impiden equiparar uniformidad con calidad. Una clase generada puede ser
estructuralmente correcta y, aun así, carecer de comportamiento funcional. Del mismo
modo, un componente especializado puede ser adecuado aunque no herede de una base
genérica.

\section{Evidencia prevista}

La evaluación combinará cuatro clases de evidencia:

\begin{itemize}[leftmargin=*]
  \item \textbf{estructural:} dependencias, contratos, metadatos y cobertura;
  \item \textbf{funcional:} resultados de operaciones y pruebas de invariantes;
  \item \textbf{experimental:} latencia, bytes, escrituras, almacenamiento y tiempo de
  reconstrucción;
  \item \textbf{analítica:} linaje, repetibilidad, corrección temporal y calidad de los
  datos derivados.
\end{itemize}

La presencia de archivos o documentación se considerará evidencia estructural, no prueba
de funcionamiento. Una capacidad solo se declarará validada si existe una ejecución
observable bajo condiciones documentadas.

\chapter{Arquitectura COBRIA}
\label{ch:arquitectura-cobria}

\section{Definición}

\begin{quote}
COBRIA es una arquitectura dirigida por metadatos que representa el estado autorizado
mediante entidades canónicas, separa el modelo semántico de su representación física y
produce proyecciones reconstruibles para operación, integración, minería de datos e
inteligencia artificial dentro de un ámbito explícito y auditable.
\end{quote}

La definición contiene cinco compromisos. Primero, existe una fuente de autoridad
identificable. Segundo, almacenar no equivale a modelar: la representación física puede
cambiar sin alterar el significado. Tercero, los datos derivados poseen contrato y no
son copias accidentales. Cuarto, la pertenencia a un organización o contexto forma parte del
diseño. Quinto, analítica e IA consumen proyecciones trazables en lugar de exportaciones
opacas.

\section{Capas y dirección de dependencias}

\begin{figure}[H]
\centering
\resizebox{0.92\textwidth}{!}{%
\begin{tikzpicture}[
  layer/.style={draw=black,minimum width=12.8cm,minimum height=1cm,align=center},
  side/.style={draw=black,dashed,minimum width=3.2cm,minimum height=1cm,align=center},
  arr/.style={-{Latex[length=2.4mm]},semithick},
  node distance=5mm and 8mm
]
\node[layer] (present) {Presentación, API y consumidores};
\node[layer,below=of present] (use) {Casos de uso: intención, autorización e idempotencia};
\node[layer,below=of use] (service) {Servicios de dominio y aplicación};
\node[layer,below=of service] (repo) {Repositorios / Mapeadores de datos};
\node[layer,below=of repo,fill=black!8] (canonical) {Entidades canónicas};
\node[layer,below=of canonical] (storage) {Pasarela de almacenamiento y motor NoSQL};
\node[side,right=of use] (meta) {COBRIA Metadatos};
\node[side,right=of canonical] (proj) {Gestores de proyecciones};
\node[side,left=of canonical] (gov) {Auditoría y gobierno};
\draw[arr] (present)--(use);
\draw[arr] (use)--(service);
\draw[arr] (service)--(repo);
\draw[arr] (repo)--(canonical);
\draw[arr] (canonical)--(storage);
\draw[dashed] (meta)--(use);
\draw[dashed] (meta)--(proj);
\draw[dashed] (proj)--(canonical);
\draw[dashed] (gov)--(canonical);
\end{tikzpicture}%
}
\caption{Capas y preocupaciones transversales de COBRIA.}
\label{fig:capas-cobria}
\end{figure}

La dirección descendente de la \cref{fig:capas-cobria} representa conocimiento de
abstracciones, no necesariamente el orden temporal de cada llamada. Un escucha del
almacenamiento puede iniciar una notificación ascendente, pero su dato debe cruzar el
Mapeador de datos antes de presentarse como entidad. La interfaz no debe construir rutas
físicas ni actualizar proyecciones por cuenta propia.

\section{Módulos internos}

\subsection{COBRIA Núcleo}

Incluye entidades, repositorios, servicios, casos de uso, catálogos y contexto. Es el
subconjunto mínimo. Un sistema pequeño puede adoptar Núcleo sin Analítica o Inteligencia.

\subsection{COBRIA Metadatos}

Contiene descripciones de campos, tipos, mapeos, ámbito, relaciones, sensibilidad,
índices y ciclo de vida. Los metadatos pueden alimentar validadores y generadores, pero
no sustituyen pruebas funcionales.

\subsection{COBRIA Proyecciones}

Define funciones de derivación, destinos, niveles de consistencia, verificadores y
procedimientos de reconstrucción. La proyección puede residir en el mismo motor o en uno
diferente.

\subsection{COBRIA Analítica}

Transforma datos autorizados en eventos, agregaciones, series, grafos, conjuntos de datos y
características versionadas. Su diseño debe preservar linaje y tiempo de observación.

\subsection{COBRIA Inteligencia}

Administra versiones de modelos, inferencias, evidencia, retroalimentación, consentimiento y
agentes. Ninguna predicción adquiere por sí sola autoridad sobre el estado canónico.

\subsection{COBRIA Gobierno}

Reúne auditoría, retención, protección de datos, observabilidad, reconciliación y
políticas de acceso. La gobernanza no se considera un módulo opcional cuando se manejan
datos sensibles o decisiones asistidas por IA.

\section{COBRIA frente a patrones relacionados}

COBRIA puede usar separación de lectura y escritura, pero no exige dos bases ni
mensajería. Por ello no es sinónimo de CQRS. CQRS permite optimizar modelos de lectura y
escritura de forma independiente, aunque añade sincronización y consistencia eventual
\parencite{microsoft2025cqrs}. COBRIA adopta esa separación cuando existe una razón
medible y añade el contrato de reconstrucción y el linaje analítico.

Tampoco exige abastecimiento por eventos. La fuente canónica puede ser el estado actual versionado,
un historial de eventos o una combinación. En un sistema basado en estado, la auditoría
no convierte automáticamente cada cambio en un evento capaz de reconstruir todo el
agregado.

Las proyecciones COBRIA incluyen vistas materializadas, pero su alcance es mayor: pueden
ser índices booleanos, relaciones inversas, resúmenes, artefactos multimedia o conjuntos de datos.
Comparten la propiedad esencial de ser desechables y regenerables desde una fuente
\parencite{microsoft2025materialized,gupta1995materialized}.

\chapter{Modelo formal de datos canónicos}
\label{ch:modelo-formal}

\section{Conjuntos fundamentales}

Sea $\mathcal{E}$ el conjunto de entidades canónicas, $\mathcal{P}$ el conjunto de
proyecciones, $\mathcal{S}$ el conjunto de ámbitos y $\mathcal{M}$ el conjunto de
versiones de metadatos. Una entidad se define como:

\begin{equation}
E=(I,A,S,V,L,R),
\label{eq:entidad}
\end{equation}

donde $I$ es la identidad, $A$ los atributos semánticos, $S\in\mathcal{S}$ el ámbito,
$V$ la versión del esquema, $L$ el estado de ciclo de vida y $R$ la revisión de
concurrencia.

La identidad global lógica puede expresarse como:

\begin{equation}
I_g=(\tau,S,I_l),
\end{equation}

donde $\tau$ es el tipo de entidad e $I_l$ el identificador local. Dos entidades con el
mismo identificador local pero ámbitos distintos no son la misma entidad.

\section{Atributos y ausencia de valor}

Cada atributo $a_j\in A$ declara tipo, nulabilidad, valor inicial, sensibilidad y reglas
de validación. COBRIA distingue entre ausencia, valor nulo y valor vacío cuando el
dominio lo requiere. Convertir indiscriminadamente los tres estados puede destruir
información necesaria para migraciones o analítica.

Los valores monetarios deben conservar moneda y unidad menor; las fechas de negocio no
deben confundirse con instantes; y los identificadores externos deben diferenciarse de
las claves internas. Estas decisiones pertenecen al modelo semántico, no al nombre corto
utilizado en almacenamiento.

\section{Ámbito y multiinquilinato}

Un ámbito puede contener varios niveles:

\begin{equation}
S=(\mathrm{organizacion},\mathrm{contexto}_1,\ldots,\mathrm{contexto}_q).
\end{equation}

Para toda operación autorizada $o$ realizada por un actor $u$ sobre una entidad $E$, se
exige:

\begin{equation}
\operatorname{autorizado}(u,o,E)\Rightarrow \operatorname{ambito}(E)\in
\operatorname{ambitos}(u,o).
\label{eq:ambito}
\end{equation}

La condición de la \cref{eq:ambito} debe verificarse en una frontera confiable. Ocultar
una entidad en la interfaz no constituye aislamiento. La ruta física puede incluir el
ámbito para acotar consultas, pero la seguridad no debe depender únicamente de que el
cliente construya correctamente esa ruta.

\section{Ciclo de vida}

El estado $L$ se modela como:

\begin{equation}
L\in\{active,archived,deleted\},
\end{equation}

con extensiones específicas del dominio. \emph{Eliminado} puede representar una marca
lógica cuando existen obligaciones de auditoría. Una política de retención define si y
cuándo procede la eliminación física. Las proyecciones deben excluir o representar
entidades archivadas de acuerdo con su contrato, no según una convención implícita.

\section{Revisión y concurrencia}

La revisión $R\in\mathbb{N}_0$ aumenta con cada modificación aceptada. Una actualización
optimista exige:

\begin{equation}
R_{request}=R_{stored}.
\end{equation}

Si la igualdad no se cumple, el caso de uso rechaza, reintenta desde un estado nuevo o
ejecuta una reconciliación explícita. El control de revisión no reemplaza transacciones
atómicas cuando una invariante involucra varias ubicaciones.

\section{Invariantes canónicas}

Se proponen cinco invariantes generales:

\begin{align}
&\forall E\in\mathcal{E}: I(E)\neq\varnothing,\\
&\forall E\in\mathcal{E}: S(E)\neq\varnothing\quad\text{si el tipo es ámbitod},\\
&\forall E_1,E_2: I_g(E_1)=I_g(E_2)\Rightarrow E_1=E_2,\\
&\forall E: deserialize(serialize(E,M),M)=E,\\
&\forall p\in\mathcal{P}: source(p)\subseteq\mathcal{E}\cup\mathcal{P}.
\end{align}

La cuarta igualdad se entiende sobre campos persistidos y admite normalizaciones
declaradas. Un timestamp generado por servidor, por ejemplo, debe modelarse como tal en
vez de ocultarse como una diferencia accidental.

\chapter{Representación física y persistencia}
\label{ch:persistencia}

\section{Mapeo semántico--físico}

Sea $M_v$ un mapeo versionado y $F$ una función de serialización:

\begin{equation}
F_{M_v}:\mathcal{E}\rightarrow\mathcal{R},
\end{equation}

donde $\mathcal{R}$ es el conjunto de representaciones físicas. La operación inversa
$G_{M_v}$ debe satisfacer:

\begin{equation}
G_{M_v}(F_{M_v}(E))\equiv E.
\label{eq:roundtrip}
\end{equation}

La equivalencia de la \cref{eq:roundtrip} incluye identidad y semántica, no propiedades
transitorias de memoria. Un ejemplo neutral es:

\begin{table}[H]
\centering
\caption{Ejemplo de mapeo semántico--físico versionado.}
\label{tab:fieldmap}
\begin{tabular}{lll}
\toprule
Campo semántico & Clave física & Regla \\
\midrule
organizaciónId & a & cadena no vacía \\
contextId & b & cadena no vacía \\
status & c & código de catálogo versionado \\
createdAt & d & instante UTC \\
revision & e & entero no negativo \\
\bottomrule
\end{tabular}
\end{table}

Las claves compactas pueden reducir repetición, pero disminuyen legibilidad y elevan el
costo de diagnóstico. COBRIA no las prescribe; exige medir su beneficio y conservar un
diccionario versionado. Una clave física nunca debe reutilizarse con otro significado
dentro de la misma versión.

\section{Repositorio y Mapeador de datos}

El Repositorio ofrece una colección conceptual de entidades; el Mapeador de datos traduce
entre objeto y almacenamiento sin introducir dependencias físicas en el dominio
\parencite{fowler2002poeaa}. En una implementación compacta ambas responsabilidades
pueden residir en una clase, siempre que su interfaz permanezca explícita.

El contrato mínimo es:

\begin{verbatim}
get(id, ámbito) -> Entidad | null
list(query, ámbito, page) -> Page<Entidad>
save(entity, expectedRevision) -> Entidad
remove(id, ámbito, policy) -> Result
subscribe(id, ámbito, observer) -> Unsubscribe
\end{verbatim}

El Repositorio no decide por sí solo una transición de negocio. Puede validar forma y
ámbito, pero una operación como aprobar, cancelar o transferir pertenece al caso de uso.

\section{Mapa de identidad y caché}

Una caché $C:I_g\rightarrow E$ evita cargar repetidamente la misma identidad durante una
unidad de trabajo. Su acceso promedio puede aproximarse a $O(1)$ mediante tabla hash,
pero introduce expiración, invalidación y riesgo de servir datos obsoletos. COBRIA exige
declarar el alcance de la caché:

\begin{itemize}[leftmargin=*]
  \item por operación, con vida corta y consistencia sencilla;
  \item por sesión, con invalidación mediante eventos o TTL;
  \item compartida, con políticas explícitas de coherencia y aislamiento.
\end{itemize}

Una caché nunca debe omitir la dimensión de organización en su clave.

\section{Pasarela de almacenamiento}

El Pasarela adapta primitivas como lectura, transacción, actualización multipath,
escucha, paginación y generación de identidad. Esta frontera permite contrastar el
modelo en emuladores y sustituir el motor, pero no promete portabilidad automática: las
garantías varían entre almacenes. Sistemas como Dynamo priorizan disponibilidad y
resolución de versiones \parencite{decandia2007dynamo}, mientras Spanner demuestra un
modelo distribuido con transacciones externamente consistentes bajo otras decisiones de
infraestructura \parencite{corbett2012spanner}.

\section{Registro de rutas semánticas}

Cuando el almacenamiento utiliza nombres físicos compactos, el desacoplamiento no se
logra solo con un Mapeador. También se necesita un registro central que traduzca nombres
semánticos a nodos y construya rutas con ámbito. Sea
\(\rho_v:S\times K^*\rightarrow Path\) la función versionada que recibe un símbolo
semántico \(S\) y segmentos de clave \(K\). Los servicios de aplicación dependen de
\(\rho_v\), no de literales físicos dispersos.

Este registro aporta cuatro propiedades comprobables:

\begin{enumerate}[leftmargin=*]
  \item una ruta física cambia en un punto controlado;
  \item una búsqueda estática puede prohibir literales fuera de infraestructura;
  \item las pruebas verifican normalización y presencia obligatoria de ámbito;
  \item telemetría, reglas y migraciones comparten el mismo vocabulario.
\end{enumerate}

El registro no convierte automáticamente un esquema compacto en modelo canónico. Su
función es establecer una frontera semántico--física y hacer localizable la deuda.

\section{Capa de servicios y casos de uso}

El Capa de servicios establece capacidades y coordina la respuesta de la aplicación
\parencite{fowler2002poeaa}. COBRIA distingue:

\begin{itemize}[leftmargin=*]
  \item \textbf{servicio de dominio:} regla que no pertenece naturalmente a una sola
  entidad;
  \item \textbf{servicio de aplicación:} composición técnica de repositorios, políticas
  y proveedores;
  \item \textbf{caso de uso:} intención completa iniciada por un actor o proceso.
\end{itemize}

Un caso de uso recibe comando, actor, ámbito, clave de idempotencia y revisión esperada.
Produce un resultado funcional, una auditoría y los cambios derivados requeridos. La
interfaz no debe convertir un cambio de estado complejo en un CRUD genérico.

\section{Contexto autorizado como dato derivado confiable}

El ámbito solicitado por cabeceras, query string o cuerpo identifica una intención, no
demuestra autorización. COBRIA distingue \(S_{solicitado}\) de
\(S_{autorizado}\). El segundo debe derivarse en servidor a partir de identidad,
declaraciones y relación comprobada con el recurso:

\[
S_{autorizado}=\operatorname{autorizar}(actor,recurso,politica).
\]

Servicios, auditoría, facturación, FinOps y proyecciones posteriores consumen únicamente
el contexto autorizado. Recalcularlo de manera distinta en cada capa crea divergencia;
aceptarlo del cliente permite atribución falsa. Por ello, el contexto validado es una
proyección efímera de seguridad: tiene fuente, política, alcance y vida limitada, pero
no reemplaza la autorización de cada operación sensible.

\chapter{Teoría de proyecciones reconstruibles}
\label{ch:proyecciones}

\section{Definición}

Una proyección $P_i$ se define mediante una función versionada:

\begin{equation}
P_i=g_{i,v}(D_i),
\label{eq:projection}
\end{equation}

donde $D_i$ es un conjunto de entidades canónicas o proyecciones previamente
autorizadas. Para evitar ciclos no controlados, el grafo de dependencias de
reconstrucción debe ser acíclico o declarar un punto fijo y un algoritmo de convergencia.

Una copia es proyección COBRIA solo si declara:

\begin{enumerate}[leftmargin=*]
  \item fuentes y ámbito;
  \item función y versión;
  \item destino y clave;
  \item propietario de escritura;
  \item nivel de consistencia;
  \item procedimiento de reconstrucción;
  \item verificador de integridad;
  \item política de privacidad y retención.
\end{enumerate}

\section{Taxonomía}

\begin{table}[H]
\centering
\caption{Taxonomía de proyecciones COBRIA.}
\label{tab:taxonomia-proyecciones}
\begin{tabularx}{\textwidth}{p{2.5cm}p{3.3cm}X}
\toprule
Clase & Consistencia esperada & Ejemplos y tratamiento \\
\midrule
Crítica & Atómica con la operación o verificada antes de confirmar & Reserva de recurso,
saldo operativo, unicidad. Si no puede mantenerse, falla el caso de uso.\\
Operacional & Eventual acotada & Índices de consulta, resúmenes de lista, relaciones
inversas. Debe existir reconciliación.\\
Analítica & Eventual y versionada & Series, características, conjuntos de datos, agregados históricos.
Prioriza linaje y corrección temporal.\\
Efímera & Sin garantía durable & Caché, vista previa o cálculo temporal. Puede descartarse sin
reconstrucción persistente.\\
\bottomrule
\end{tabularx}
\end{table}

Clasificar no convierte una proyección en correcta. La clase define qué retrasos y
fallos son aceptables y qué pruebas se requieren.

\section{Propiedades}

\subsection{Determinismo}

Con fuentes y versión iguales, una proyección determinista produce el mismo resultado:

\begin{equation}
D_i^{(1)}=D_i^{(2)}\Rightarrow g_{i,v}(D_i^{(1)})=g_{i,v}(D_i^{(2)}).
\end{equation}

Si utiliza tiempo actual, aleatoriedad o un proveedor externo, esos valores deben
capturarse como entradas versionadas.

\subsection{Idempotencia}

Aplicar dos veces el mismo cambio no debe duplicar resultados:

\begin{equation}
apply(apply(P,c),c)=apply(P,c).
\end{equation}

\subsection{Reconstructibilidad}

Si $P_i$ se pierde, debe poder calcularse $P_i'$ tal que:

\begin{equation}
verify(P_i',g_{i,v}(D_i))=true.
\end{equation}

\subsection{No autoridad}

Una proyección no se edita para corregir negocio. La corrección se ejecuta sobre el
canónico mediante un caso de uso y luego se propaga.

\section{Grafo de dependencias}

Sea $G=(V,E_d)$ un grafo donde $V=\mathcal{E}\cup\mathcal{P}$ y una arista $(x,p)$
indica que $p$ depende de $x$. El orden de reconstrucción se obtiene mediante orden
topológico. Si se detecta un ciclo no declarado, el plan se rechaza.

\begin{figure}[H]
\centering
\begin{tikzpicture}[
  n/.style={draw=black,rounded corners=1pt,minimum width=3.2cm,minimum height=8mm,align=center},
  p/.style={draw=black,fill=black!7,minimum width=3.2cm,minimum height=8mm,align=center},
  a/.style={-{Latex[length=2.2mm]},semithick},node distance=9mm and 14mm]
\node[n] (e1) {Entidad canónica A};
\node[n,below=of e1] (e2) {Entidad canónica B};
\node[p,right=of e1] (idx) {Índice operacional};
\node[p,right=of e2] (sum) {Resumen compuesto};
\node[p,right=of sum] (data) {Conjunto de datos versionado};
\draw[a] (e1)--(idx);
\draw[a] (e1)--(sum);
\draw[a] (e2)--(sum);
\draw[a] (sum)--(data);
\end{tikzpicture}
\caption{Ejemplo de grafo acíclico de reconstrucción.}
\label{fig:grafo-proyecciones}
\end{figure}

\section{Contrato de proyección}

Una representación neutral puede adoptar la siguiente forma:

\begin{verbatim}
ProjectionContract {
  name, version, class,
  sources[], ámbito,
  destination, keyFunction,
  build, verify,
  consistency, responsable,
  privacy, retention
}
\end{verbatim}

El contrato es declarativo respecto al propósito, aunque `build` y `verify` pueden ser
funciones. Debe poder responder quién escribe la proyección y cómo se repara. Si esas
preguntas no tienen respuesta, la duplicación permanece fuera de gobierno.

\chapter{Algoritmos, consistencia y recuperación}
\label{ch:algoritmos-consistencia}

\section{Creación de una entidad}

\begin{verbatim}
algorithm Create(command, actor, idempotencyKey):
  require authorize(actor, command.ámbito, command.action)
  if resultExists(idempotencyKey): return previousResult
  metadata <- registry.forType(command.type)
  entity <- metadata.construct(command.carga útil)
  validate(entity)
  critical <- projections.criticalFor(entity.type)
  atomicWrite(entity, buildAll(critical, entity))
  enqueue(projections.nonCriticalFor(entity.type), entity.revision)
  audit(actor, command, entity)
  return remember(idempotencyKey, entity)
\end{verbatim}

Las proyecciones críticas forman parte de la confirmación. Las demás se envían a un
proceso idempotente. Si el motor no ofrece una escritura atómica suficiente, el diseño
debe reducir la frontera crítica o implementar una reserva y compensación explícitas.

\section{Actualización optimista}

\begin{verbatim}
algorithm Update(command, actor, expectedRevision):
  current <- repository.get(command.id, command.ámbito)
  require current.revision = expectedRevision
  next <- useCase.transition(current, command)
  next.revision <- current.revision + 1
  validateTransition(current, next)
  atomicWrite(next, criticalDelta(current, next))
  enqueue(nonCriticalDelta(current, next))
  return next
\end{verbatim}

El delta debe eliminar claves antiguas además de añadir nuevas. Cambiar el estado o el
ámbito sin retirar índices previos produce referencias fantasma.

\section{Reconstrucción total}

\begin{verbatim}
algorithm RebuildAll(contract, ámbito):
  acquireLease(contract.name, ámbito)
  targetVersion <- createShadowDestination()
  for page in sourcePages(contract.sources, ámbito):
    records <- contract.build(page)
    writeIdempotently(targetVersion, records)
    checkpoint(page.cursor)
  report <- contract.verify(targetVersion)
  require report.errors = 0
  atomicCutover(targetVersion)
  releaseLease()
  return report
\end{verbatim}

El destino sombra evita dejar una proyección parcialmente reconstruida como si estuviera
completa. Para volúmenes pequeños puede reconstruirse en sitio, pero debe exponerse el
estado `rebuilding` y tolerar reanudación.

\section{Reconstrucción incremental}

\begin{verbatim}
algorithm RebuildChanged(contract, marca de corte):
  changes <- canonicalChangesAfter(marca de corte)
  for change in ordered(changes):
    require notProcessed(change.id, contract.version)
    delta <- contract.buildDelta(change.before, change.after)
    applyIdempotently(delta)
    markProcessed(change.id, contract.version)
  return contract.verifyAffected(changes)
\end{verbatim}

La reconstrucción incremental necesita un registro confiable de cambios o una marca de
actualización consultable. Si no existe, escanear el conjunto canónico puede ser la única
forma correcta de detectar diferencias.

\section{Verificación de integridad}

Se definen cuatro clases de anomalía:

\begin{align}
orphan &= \{p\in P: sourceId(p)\notin I(\mathcal{E})\},\\
missing &= \{e\in E: expected(e)\setminus actual(e)\neq\varnothing\},\\
stale &= \{p\in P: sourceRevision(p)<revision(source(p))\},\\
divergent &= \{p\in P:p\neq g_{v}(source(p))\}.
\end{align}

Un reporte debe separar conteos, ejemplos anonimizados y severidad. La reparación
automática es apropiada para derivados deterministas; una divergencia canónica requiere
revisión de negocio.

\section{Consistencia por clase}

\begin{table}[H]
\centering
\caption{Mecanismos de consistencia recomendados.}
\label{tab:consistencia-mecanismos}
\begin{tabularx}{\textwidth}{p{3cm}p{4cm}X}
\toprule
Necesidad & Mecanismo & Consideración \\
\midrule
Invariante local & Transacción o compare-and-set & Rechazar revisiones obsoletas.\\
Varias rutas del mismo motor & Actualización multipath atómica & Validar límites reales
del proveedor.\\
Derivado operacional & Evento/cola de salida y proceso de trabajo idempotente & Medir retraso y reconciliar.\\
Integración externa & Saga o compensación & No asumir transacción distribuida.\\
Conjunto de datos analítico & Lote versionado con marca de corte & Preservar tiempo y linaje.\\
\bottomrule
\end{tabularx}
\end{table}

La consistencia fuerte global no debe afirmarse por analogía. Las garantías dependen del
motor, partición y frontera transaccional. COBRIA documenta esas garantías en el contrato
del caso de uso.

\section{Estado materializado: completitud, frescura y procedencia}

Una vista no debe considerarse válida por el solo hecho de existir. El contrato mínimo
incluye metadatos de control:

\begin{verbatim}
_meta:
  complete: boolean
  updatedAt: timestamp
  schemaVersion: integer
  sourceRevision|marca de corte: value
  count|checksum: optional
\end{verbatim}

La lectura acepta la vista únicamente si está completa, dentro de su presupuesto de
frescura y en una versión compatible. Una mutación puede actualizarla, invalidarla o
marcarla incompleta. Si la proyección es mejor esfuerzo, su fallo no revierte la escritura
canónica; la lectura cae a la fuente autorizada y programa reparación.

Cuando varias solicitudes encuentran una vista fría, la reconstrucción debe
\emph{coalescerse}: una promesa, arrendamiento o bloqueo por ámbito evita que todas repitan el mismo
escaneo. Esta protección contra \emph{cache stampede} forma parte del contrato de
reconstrucción y debe probarse con solicitudes concurrentes.

\section{Arquitectura ejecutable y funciones de aptitud}

Una regla escrita puede degradarse silenciosamente. COBRIA incorpora funciones de
aptitud arquitectónica: pruebas o análisis estáticos que convierten una propiedad
estructural en una condición de integración continua. Ejemplos:

\begin{itemize}[leftmargin=*]
  \item presentación no importa infraestructura de datos;
  \item rutas físicas compactas no aparecen en servicios de aplicación;
  \item todo modelo de lectura declara metadatos de completitud y frescura;
  \item una operación sensible usa contexto autorizado por servidor;
  \item el número de consumidores de una fachada heredado no puede crecer;
  \item rutas críticas respetan presupuestos de bytes, LCP y bloqueo.
\end{itemize}

Para una deuda medible \(d_t\), el control automatizado de migración impone
\(d_{t+1}\leq d_t\), y el objetivo final puede fijarse en cero. Este presupuesto
monótono permite migrar sin una reescritura total y evita que una fachada temporal se
convierta en arquitectura permanente.

\section{Cola de salida, reintentos y duplicados}

Cuando guardar el canónico y publicar un mensaje no forman una transacción única, existe
una ventana de pérdida. El patrón cola de salida conserva la intención de publicación junto al
cambio autorizado; un proceso de trabajo la entrega y registra el resultado. La entrega puede ser
al menos una vez, por lo que consumidores y Gestores de proyecciones deben ser idempotentes.

Una clave de idempotencia se asocia con actor, ámbito, operación y ventana temporal. No
debe reutilizarse entre organizaciones. El resultado previo puede devolverse siempre que los
parámetros relevantes coincidan; en caso contrario, la reutilización se rechaza.

\section{Recuperación y observabilidad}

Cada proceso de reconstrucción registra:

\begin{itemize}[leftmargin=*]
  \item contrato y versión;
  \item ámbito y rango temporal;
  \item cursor o checkpoint;
  \item entidades leídas y registros escritos;
  \item anomalías antes y después;
  \item duración y consumo aproximado;
  \item identidad del operador o proceso de trabajo;
  \item resultado del cambio de versión.
\end{itemize}

El objetivo de recuperación (RTO) y el punto de recuperación (RPO) deben definirse por
clase de proyección. Una caché puede perderse sin RPO; un saldo derivado utilizado en una
operación exige un tratamiento diferente.

\section{Antipatrones}

COBRIA identifica como antipatrones:

\begin{enumerate}[leftmargin=*]
  \item editar una proyección para corregir la fuente;
  \item construir rutas físicas en la presentación;
  \item mantener el mismo índice desde múltiples servicios sin propietario;
  \item usar un ID sin ámbito como clave de caché multiinquilino;
  \item generar una clase por cada estructura sin justificar identidad o ciclo de vida;
  \item llamar transacción a una secuencia de escrituras independientes;
  \item entrenar con una exportación sin versión, marca de corte ni linaje;
  \item permitir que un agente de IA escriba directamente en el almacenamiento;
  \item declarar una integración completa por la sola presencia de modelos o contratos;
  \item añadir proyecciones sin medir amplificación y frecuencia de consulta;
  \item confiar en el ámbito enviado por el cliente para auditoría o autorización;
  \item servir una vista existente sin comprobar completitud, versión y frescura;
  \item crear una capa de compatibilidad sin inventario ni presupuesto decreciente;
  \item documentar límites de capas sin un control automatizado que detecte regresiones.
\end{enumerate}

\section{Criterio mínimo de adopción}

COBRIA Núcleo puede justificarse cuando existe al menos una de estas condiciones: dominio
con reglas relevantes, multiinquilinato, duplicación derivada, necesidad de reconstrucción,
procesamiento asíncrono o uso analítico trazable. Si el sistema es un CRUD pequeño, con
pocos registros y consultas directas, Repositorio más un modelo sencillo puede ser
suficiente. La arquitectura se considera exitosa cuando hace explícitos los compromisos,
incluso si la decisión resultante es no crear una proyección.

\section{Síntesis de la Parte II}

Esta parte transformó la intuición inicial en un artefacto definido. COBRIA organiza el
estado mediante entidades con identidad, ámbito, versión, ciclo de vida y revisión;
separa semántica de persistencia mediante un mapeo reversible; y asigna a Repositorio,
Servicio y Caso de uso fronteras distintas. Su elemento diferenciador es el contrato de
proyección, que convierte índices, resúmenes y conjuntos de datos en derivaciones gobernadas y
reconstruibles.

También se establecieron algoritmos de creación, actualización, reconstrucción y
verificación, junto con mecanismos de consistencia proporcionales a la criticidad. Estas
definiciones preparan la siguiente parte del estudio: analizar complejidad, amplificación
y costo, y comprobar mediante experimentos si la disciplina propuesta produce un
beneficio neto.
